% SOURCE: original Chinese PDF, viewer pages 19--24.
\section{Position and Momentum Representations}
\subsection*{1) Continuous spectrum}
Unlike the spin angular momentum $\op s_z$, which takes discrete values, position $\op x$ and momentum $\op p$ take continuous values.

The eigenstates of position or momentum constitute a continuous position or momentum representation.

Eigenvalue equations:
\begin{equation}
\op x\ket x=x\ket x,\qquad\op p\ket p=p\ket p.
\end{equation}
Basis vectors: $\ket x$, $\ket p$.

Orthonormality:
\begin{equation}
\braket x{x'}=\delta(x-x'),\qquad\braket p{p'}=\delta(p-p').
\end{equation}
Completeness:
\begin{equation}
\int\dd x\,\ket x\bra x=\mathbf1,\qquad
\int\dd p\,\ket p\bra p=\mathbf1.
\end{equation}
For an arbitrary state $\ket\psi$,
\begin{equation}
\ket\psi=\int\dd x\,\ket x\braket x\psi
=\int\dd p\,\ket p\braket p\psi.
\end{equation}
The functions $\braket x\psi$ and $\braket p\psi$ are the concrete forms of $\ket\psi$ in the position and momentum representations; they are continuous column matrices.

In three-dimensional space,
\begin{equation}
\op{\vect x}\ket{\vect x}=\vect x\ket{\vect x},\qquad
\op{\vect p}\ket{\vect p}=\vect p\ket{\vect p}.
\end{equation}
\textbf{Postulate of quantum mechanics:} The fundamental commutation relations are
\begin{equation}
[\op x_i,\op x_j]=0,\qquad [\op p_i,\op p_j]=0,\qquad
[\op x_i,\op p_j]=\ii\hbar\delta_{ij},\quad i,j=1,2,3.
\end{equation}
The position--momentum uncertainty relation is
\begin{equation}
\langle(\Delta\op x)^2\rangle\langle(\Delta\op p_x)^2\rangle
\ge\left(\frac1{2\ii}\langle[\op x,\op p_x]\rangle\right)^2=\frac{\hbar^2}{4}.
\end{equation}
\subsection*{2) Spatial translations}
Define the spatial-translation operator $\op U_{\mathrm{tr}}$, with the subscript ``tr'' indicating translation (distinct from the scattering transition operator $\op T$ introduced later), by
\begin{equation}
\op U_{\mathrm{tr}}(\dd\vect x)\ket{\vect x}=\ket{\vect x+\dd\vect x}.
\end{equation}
For an arbitrary state $\ket\psi$,
\begin{equation}\begin{aligned}
\op U_{\mathrm{tr}}(\dd\vect x)\ket\psi
&=\op U_{\mathrm{tr}}(\dd\vect x)\int\dd^3x\,\ket{\vect x}\braket{\vect x}\psi\\
&=\int\dd^3x\,\ket{\vect x+\dd\vect x}\braket{\vect x}\psi
&&\text{(the operator acts only on vectors)}\\
&=\int\dd^3x\,\ket{\vect x}\braket{\vect x-\dd\vect x}\psi
&&\text{(change of integration variable).}
\end{aligned}\end{equation}
For an infinitesimal translation, take
\begin{equation}
\op U_{\mathrm{tr}}(\dd\vect x)=\mathbf1-\ii\op{\vect K}\cdot\dd\vect x,
\end{equation}
where the generator $\op{\vect K}$ is Hermitian.

A translation is a unitary operator:
\begin{equation}
\op U_{\mathrm{tr}}(\dd\vect x)\op U_{\mathrm{tr}}^\dagger(\dd\vect x)
=(\mathbf1-\ii\op{\vect K}\cdot\dd\vect x)
 (\mathbf1+\ii\op{\vect K}\cdot\dd\vect x)
=\mathbf1,
\end{equation}
neglecting second-order infinitesimals.

Composition law for translations:
\begin{equation}\begin{aligned}
\op U_{\mathrm{tr}}(\dd\vect x_1)\op U_{\mathrm{tr}}(\dd\vect x_2)
&=(\mathbf1-\ii\op{\vect K}\cdot\dd\vect x_1)
(\mathbf1-\ii\op{\vect K}\cdot\dd\vect x_2)\\
&=\mathbf1-\ii\op{\vect K}\cdot(\dd\vect x_1+\dd\vect x_2)
=\op U_{\mathrm{tr}}(\dd\vect x_1+\dd\vect x_2).
\end{aligned}\end{equation}
A finite translation is obtained by repeated infinitesimal translations.

For an arbitrary state $\ket\psi$,
\begin{equation}\begin{aligned}
\op{\vect x}\op U_{\mathrm{tr}}(\dd\vect x)\ket\psi
&=\int\dd^3x\,\op{\vect x}\ket{\vect x+\dd\vect x}\braket{\vect x}\psi\\
&=\int\dd^3x\,(\vect x+\dd\vect x)\ket{\vect x+\dd\vect x}\braket{\vect x}\psi,\\
\op U_{\mathrm{tr}}(\dd\vect x)\op{\vect x}\ket\psi
&=\int\dd^3x\,\op U_{\mathrm{tr}}(\dd\vect x)\vect x\ket{\vect x}\braket{\vect x}\psi\\
&=\int\dd^3x\,\vect x\ket{\vect x+\dd\vect x}\braket{\vect x}\psi.
\end{aligned}\end{equation}
Hence
\begin{equation}\begin{aligned}
[\op{\vect x},\op U_{\mathrm{tr}}(\dd\vect x)]\ket\psi
&=\int\dd^3x\,\dd\vect x\,\ket{\vect x+\dd\vect x}\braket{\vect x}\psi\\
&=\int\dd^3x\,\dd\vect x\,\ket{\vect x}\braket{\vect x}\psi
=\dd\vect x\ket\psi,
\end{aligned}\end{equation}
where second-order infinitesimals have been neglected. Thus
\begin{equation}
[\op{\vect x},\op U_{\mathrm{tr}}(\dd\vect x)]=\dd\vect x,\qquad
[\op{\vect x},\op{\vect K}\cdot\dd\vect x]=\ii\dd\vect x.
\end{equation}
If the translation is taken in the $x_j$ direction, $\dd\vect x=\epsilon\vect e_j$, then
\begin{equation}
[\op x_i,\op K_j]=\ii\delta_{ij}.
\end{equation}
Comparing with the fundamental quantum-mechanical commutation relation,
\begin{equation}
[\op x_i,\op p_j]=\ii\hbar\delta_{ij},
\end{equation}
we have
\begin{equation}
\op{\vect p}=\hbar\op{\vect K}.
\end{equation}
Momentum is therefore the generator of translations, and
\begin{equation}
\boxed{\op U_{\mathrm{tr}}(\dd\vect x)=\mathbf1-\frac\ii\hbar\op{\vect p}\cdot\dd\vect x.}
\end{equation}
Since $[\op p_x,\op p_y]=0$,
\begin{equation}
[\op U_{\mathrm{tr}}(\dd x),\op U_{\mathrm{tr}}(\dd y)]
=\left[\mathbf1-\frac\ii\hbar\op p_x\dd x,
        \mathbf1-\frac\ii\hbar\op p_y\dd y\right]=0.
\end{equation}
This shows that the order of translations is immaterial. Translating $A\to C$ by $\dd x$ and then $C\to B$ by $\dd y$ gives the same result as translating $A\to D$ by $\dd y$ and then $D\to B$ by $\dd x$.
\begin{center}\begin{tikzpicture}[>=Latex,scale=.85]
\draw[->](-.5,0)--(4.3,0)node[right]{$x$};\draw[->](0,-.5)--(0,3.3)node[above]{$y$};
\coordinate(A)at(1,0.45);\coordinate(C)at(3.1,.45);\coordinate(D)at(1,2.5);\coordinate(B)at(3.1,2.5);
\draw[->,thick](A)--(C)node[midway,below,yshift=-4pt]{$\dd x$};\draw[->,thick](C)--(B)node[midway,right,xshift=3pt]{$\dd y$};
\draw[->,thick](A)--(D);\draw[->,thick](D)--(B);
\foreach \name in {A,B,C,D}{\node[above left]at(\name){$\name$};}
\end{tikzpicture}\end{center}
\subsection*{3) Canonical quantization}
Dirac, 1925: quantization of classical mechanical systems.

The Poisson bracket of analytical mechanics is replaced by a commutation relation in quantum mechanics:
\begin{equation}
[A,B]_{\rm P}=\sum_s\left(
\frac{\partial A}{\partial q_s}\frac{\partial B}{\partial p_s}
-\frac{\partial A}{\partial p_s}\frac{\partial B}{\partial q_s}
\right)
\longrightarrow\frac1{\ii\hbar}[\op A,\op B]
=\frac1{\ii\hbar}(\op A\op B-\op B\op A).
\end{equation}
For example,
\begin{equation}
[x_i,p_j]_{\rm P}=\delta_{ij}
\longrightarrow [\op x_i,\op p_j]=\ii\hbar\delta_{ij}.
\end{equation}
\subsection*{4) Position representation}
The matrix elements of $\op{\vect x}$ are
\begin{equation}
\mel{\vect x_1}{\op{\vect x}}{\vect x_2}
=\vect x_2\delta^{(3)}(\vect x_1-\vect x_2)
=\vect x_1\delta^{(3)}(\vect x_1-\vect x_2),
\end{equation}
so the position operator is diagonal.

From
\begin{equation}
\op U_{\mathrm{tr}}(\dd\vect x)\ket\beta
=\left(\mathbf1-\frac\ii\hbar\op{\vect p}\cdot\dd\vect x\right)\ket\beta
\end{equation}
and
\begin{equation}\begin{aligned}
\op U_{\mathrm{tr}}(\dd\vect x)\ket\beta
&=\int\dd^3x\,\ket{\vect x}\braket{\vect x-\dd\vect x}\beta\\
&=\int\dd^3x\,\ket{\vect x}\left(
\braket{\vect x}\beta-\dd\vect x\cdot\vect\nabla_{\vect x}\braket{\vect x}\beta
\right)\\
&=\ket\beta-\dd\vect x\cdot\int\dd^3x\,\ket{\vect x}\vect\nabla_{\vect x}\braket{\vect x}\beta,
\end{aligned}\end{equation}
we obtain
\begin{equation}
\op{\vect p}\ket\beta
=\int\dd^3x\,\ket{\vect x}(-\ii\hbar\vect\nabla_{\vect x})\braket{\vect x}\beta,
\end{equation}
and
\begin{equation}
\mel\alpha{\op{\vect p}}\beta
=\int\dd^3x\,\braket\alpha{\vect x}(-\ii\hbar\vect\nabla_{\vect x})\braket{\vect x}\beta.
\end{equation}
Taking $\ket\alpha=\ket{\vect x_1}$ and $\ket\beta=\ket{\vect x_2}$ gives the matrix elements of momentum in the position representation:
\begin{equation}
\mel{\vect x_1}{\op{\vect p}}{\vect x_2}
=\int\dd^3x\,\braket{\vect x_1}{\vect x}(-\ii\hbar\vect\nabla_{\vect x})\braket{\vect x}{\vect x_2}
=-\ii\hbar\vect\nabla_{\vect x_1}\delta^{(3)}(\vect x_1-\vect x_2).
\end{equation}
Now take $\ket\alpha=\ket{\vect x_1}$ and $\ket\beta=\ket{\vect p}$:
\begin{equation}
\mel{\vect x_1}{\op{\vect p}}{\vect p}
=\int\dd^3x\,\braket{\vect x_1}{\vect x}(-\ii\hbar\vect\nabla_{\vect x})\braket{\vect x}{\vect p}
=-\ii\hbar\vect\nabla_{\vect x_1}\braket{\vect x_1}{\vect p}.
\end{equation}
But the momentum eigenvalue equation gives
\begin{equation}
\vect p\braket{\vect x_1}{\vect p}
=-\ii\hbar\vect\nabla_{\vect x_1}\braket{\vect x_1}{\vect p}.
\end{equation}
Solving this first-order differential equation, we find that the momentum eigenstate in the position representation is a plane wave:
\begin{equation}
\braket{\vect x}{\vect p}=N\exp\left(\frac\ii\hbar\vect p\cdot\vect x\right).
\end{equation}
Using
\begin{equation}\begin{aligned}
\braket{\vect x_1}{\vect x_2}
&=\int\dd^3p\,\braket{\vect x_1}{\vect p}\braket{\vect p}{\vect x_2}\\
&=|N|^2\int\dd^3p\,\exp\left[\frac\ii\hbar\vect p\cdot(\vect x_1-\vect x_2)\right],
\end{aligned}\end{equation}
and
\begin{equation}
\delta^{(3)}(\vect x_1-\vect x_2)
=|N|^2(2\pi\hbar)^3\delta^{(3)}(\vect x_1-\vect x_2),
\end{equation}
we obtain
\begin{equation}
N=\frac1{(2\pi\hbar)^{3/2}},\qquad
\boxed{\braket{\vect x}{\vect p}
=\frac1{(2\pi\hbar)^{3/2}}\exp\left(\frac\ii\hbar\vect p\cdot\vect x\right).}
\end{equation}
\subsection*{5) Momentum representation}
The matrix elements of the momentum operator are
\begin{equation}
\mel{\vect p_1}{\op{\vect p}}{\vect p_2}
=\vect p_1\delta^{(3)}(\vect p_1-\vect p_2),
\end{equation}
so the momentum operator is diagonal.

For the position operator,
\begin{equation}\begin{aligned}
\mel{\vect p_1}{\op{\vect x}}{\vect p_2}
&=\int\dd^3x_1\dd^3x_2\,
\braket{\vect p_1}{\vect x_1}
\mel{\vect x_1}{\op{\vect x}}{\vect x_2}
\braket{\vect x_2}{\vect p_2}\\
&=\frac1{(2\pi\hbar)^3}\int\dd^3x_1\dd^3x_2\,
\exp\left(-\frac\ii\hbar\vect p_1\cdot\vect x_1\right)
\vect x_2\delta^{(3)}(\vect x_1-\vect x_2)
\exp\left(\frac\ii\hbar\vect p_2\cdot\vect x_2\right)\\
&=\frac1{(2\pi\hbar)^3}\int\dd^3x_1\,
\vect x_1\exp\left[\frac\ii\hbar(\vect p_2-\vect p_1)\cdot\vect x_1\right]\\
&=\ii\hbar\vect\nabla_{\vect p_1}\,
\frac1{(2\pi\hbar)^3}\int\dd^3x_1\,
\exp\left[\frac\ii\hbar(\vect p_2-\vect p_1)\cdot\vect x_1\right]\\
&=\ii\hbar\vect\nabla_{\vect p_1}\delta^{(3)}(\vect p_1-\vect p_2).
\end{aligned}\end{equation}
Note that differentiation here is with respect to momentum.

The form of the position eigenstate in the momentum representation is
\begin{equation}
\boxed{\braket{\vect p}{\vect x}
=\braket{\vect x}{\vect p}^*
=\frac1{(2\pi\hbar)^{3/2}}\exp\left(-\frac\ii\hbar\vect p\cdot\vect x\right).}
\end{equation}
\paragraph{Example 1: Prove that $\op x$ and $\op p$ are Hermitian operators in the position representation.}
For position,
\begin{equation}
\mel{x_1}{\op x}{x_2}=\delta(x_1-x_2)x_1,
\end{equation}
and
\begin{equation}\begin{aligned}
\mel{x_1}{\op x^\dagger}{x_2}
&=\mel{x_2}{\op x}{x_1}^*
=[\delta(x_2-x_1)x_2]^*
=\delta(x_2-x_1)x_2\\
&=\delta(x_1-x_2)x_1
=\mel{x_1}{\op x}{x_2}.
\end{aligned}\end{equation}
For momentum,
\begin{equation}
\mel{x_1}{\op p}{x_2}
=-\ii\hbar\frac{\partial}{\partial x_1}\delta(x_1-x_2)
=-\ii\hbar\frac{\partial}{\partial(x_1-x_2)}\delta(x_1-x_2).
\end{equation}
Also,
\begin{equation}\begin{aligned}
\mel{x_1}{\op p^\dagger}{x_2}
&=\mel{x_2}{\op p}{x_1}^*
=\left[-\ii\hbar\frac{\partial}{\partial(x_2-x_1)}\delta(x_2-x_1)\right]^*\\
&=\ii\hbar\frac{\partial}{\partial(x_2-x_1)}\delta(x_2-x_1)
=-\ii\hbar\frac{\partial}{\partial(x_1-x_2)}\delta(x_1-x_2)\\
&=\mel{x_1}{\op p}{x_2}.
\end{aligned}\end{equation}
Thus both $\op x$ and $\op p$ are Hermitian.
\paragraph{Example 2: Compute $[\op x,\op p]$ in the position representation.}
\begin{equation}\begin{aligned}
\op x\op p
&=\int\dd x_1\dd x_2\,
\op x\ket{x_1}\mel{x_1}{\op p}{x_2}\bra{x_2}\\
&=\int\dd x_1\dd x_2\,\ket{x_1}\,x_1
\left(-\ii\hbar\frac\partial{\partial x_1}\right)\delta(x_1-x_2)\bra{x_2}\\
&=-\int\dd x_1\dd x_2\,\ket{x_1}
\left(-\ii\hbar\frac\partial{\partial x_1}\right)[x_1\bra{x_1}]\,\delta(x_1-x_2)
\quad\text{(integration by parts)}\\
&=-\int\dd x_1\,\ket{x_1}
\left(-\ii\hbar\frac\partial{\partial x_1}\right)[x_1\bra{x_1}]\\
&=\int\dd x_1\,x_1\ket{x_1}
\left(-\ii\hbar\frac\partial{\partial x_1}\right)\bra{x_1}
\quad\text{(integration by parts).}
\end{aligned}\end{equation}
Similarly,
\begin{equation}
\op p\op x=-\ii\hbar+\int\dd x_1\,x_1\ket{x_1}
\left(-\ii\hbar\frac\partial{\partial x_1}\right)\bra{x_1}.
\end{equation}
Therefore
\begin{equation}
[\op x,\op p]=\ii\hbar.
\end{equation}
This shows that commutation relations do not depend on the representation.

According to the uncertainty relation, position and momentum cannot simultaneously have definite values. For example, in a momentum eigenstate the momentum has a definite value, $\mel p{(\Delta\op p)^2}p=0$, but the probability for position to take value $x$ is
\begin{equation}
|\braket x p|^2=\left|\frac1{\sqrt{2\pi\hbar}}e^{\ii px/\hbar}\right|^2
=\text{const}.
\end{equation}
The particle therefore has equal probability of appearing anywhere in $-\infty<x<\infty$; position is completely indefinite, with variance $\mel p{(\Delta\op x)^2}p=\infty$.
\paragraph{Example 3: Minimum-uncertainty wave packet.}
For equality to hold in the uncertainty relation
\begin{equation}
\langle(\Delta\op A)^2\rangle\langle(\Delta\op B)^2\rangle
\ge\left(\frac1{2\ii}\langle[\op A,\op B]\rangle\right)^2,
\end{equation}
so that the uncertainty is minimal, the state must satisfy:
\begin{enumerate}[label=\arabic*)]
\item Equality in the Schwarz inequality:
\begin{equation}
\langle(\Delta\op A)^2\rangle\langle(\Delta\op B)^2\rangle
=|\langle\Delta\op A\Delta\op B\rangle|^2.
\end{equation}
\item $\operatorname{Re}\langle\Delta\op A\Delta\op B\rangle=0$.
\end{enumerate}
The solution of condition 1 is
\begin{equation}
\Delta\op B\ket\psi=c\Delta\op A\ket\psi,
\end{equation}
where $c$ is a constant. Condition 2 implies
\begin{equation}
\operatorname{Re}\left(c\langle(\Delta\op A)^2\rangle\right)=0.
\end{equation}
Since $\langle(\Delta\op A)^2\rangle$ is real, $c=\ii a$, with $a$ real. Hence
\begin{equation}
(\op B-\langle\op B\rangle)\ket\psi
=\ii a(\op A-\langle\op A\rangle)\ket\psi.
\end{equation}
This is the constraint on $\ket\psi$ for minimum uncertainty. Choose
\begin{equation}
\op A=\op x,\qquad\op B=\op p,
\end{equation}
so
\begin{equation}
(\op p-\langle\op p\rangle)\ket\psi
=\ii a(\op x-\langle\op x\rangle)\ket\psi.
\end{equation}
In the position representation,
\begin{equation}
\left(-\ii\hbar\frac\dd{\dd x}-\langle\op p\rangle\right)\psi(x)
=\ii a(x-\langle\op x\rangle)\psi(x),\qquad
\psi(x)=\braket x\psi.
\end{equation}
The solution is
\begin{equation}
\psi(x)=A\exp\left[-\frac{a(x-\langle\op x\rangle)^2}{2\hbar}\right]
\exp\left(\frac\ii\hbar\langle\op p\rangle x\right),
\end{equation}
a Gaussian wave packet in position space.
\paragraph{Example 4: The ground state of the hydrogen atom.}
Given the state in the position representation
\begin{equation}
\braket{\vect r}\psi
=\frac1{\sqrt{\pi a_0^3}}e^{-r/a_0}\qquad\text{(hydrogen-atom ground state)},
\end{equation}
find the expectation value of the momentum.

Since $\braket{\vect r}\psi$ is not a plane wave, momentum has no definite value. The probability for momentum $\vect p$ is $|\braket{\vect p}\psi|^2$, where
\begin{equation}\begin{aligned}
\braket{\vect p}\psi
&=\int\dd^3r\,\braket{\vect p}{\vect r}\braket{\vect r}\psi
=\int\dd^3r\,\braket{\vect r}{\vect p}^*\braket{\vect r}\psi\\
&=\int\dd^3r\,\frac1{(2\pi\hbar)^{3/2}}
\exp\left(-\frac\ii\hbar\vect p\cdot\vect r\right)
\frac1{\sqrt{\pi a_0^3}}e^{-r/a_0}\\
&=\frac{(2a_0\hbar)^{3/2}\hbar}{\pi(a_0^2p^2+\hbar^2)^2}.
\end{aligned}\end{equation}
Since $\braket{\vect p}\psi$ depends only on the magnitude of momentum, the expectation value of the momentum magnitude is
\begin{equation}
\langle p\rangle=\frac{\displaystyle\int\dd p\,p^3|\braket{\vect p}\psi|^2}
{\displaystyle\int\dd p\,p^2|\braket{\vect p}\psi|^2}.
\end{equation}